% Results Section — streamlined (2026-05-11)
% Over results_v3.tex (2026-04-17):
%   (1) 2-component elevated to primary narrative in R4; retinal family as comparison.
%   (2) Citation style ported from natbib to apacite (\citep→\cite, \citet→\citeA,
%       \citealp→\citeNP).
%   (3) Sub-09 ΔRDM attribution corrected: p=0.026 = R+C (Δλ=19.5nm, g=−1.10),
%       NOT Machado Δλ=13.5nm. Headline shifted to 2-comp V1 p=0.007.
%   (4) Emery numerical β_s convergence claim removed (grounding = model structure,
%       not parameter value; per feedback_physiological_grounding.md).
%   (5) Behavioral concordance ¶ added to R3 (LOCO 100% vs SRM 33%).
%   (6) "2-comp dominated Machado" phrasing avoided; framed on bijectivity +
%       unified parameterization.

\section{Results}
\label{sec:results}

%-----------------------------------------------------------------------
\subsection{Clinical phenotype and participants}
\label{sec:results:behavior}

Two CVD participants with contrasting severity and diagnostic axis were
recruited alongside seven healthy controls (HC). Both CVD participants were
male and completed the 14-plate Ishihara pseudoisochromatic
test~\cite{ishihara1917}. Sub-08 was classified as mild--moderate
deuteranomalous (8/14 plates misread; protan--deutan discrimination plate
indicating deutan axis); Sub-09 as moderate--severe protanomalous (11/14
plates misread; protan axis). These two cases span the deutan vs.\ protan
dichotomy and the mild vs.\ moderate severity range that a clinically useful
filter must accommodate.

Same-day 8-alternative forced-choice identification of the eight scanner
stimuli confirmed the Ishihara categorisation (HC accuracy
$0.94 \pm 0.04$, Sub-09 $= 0.62$, Sub-08 $= 0.71$; chance $= 0.125$).
Sub-08's errors clustered on the red--green confusion axis (colours c2/c3
vs.\ c7; mean confusion $28\%$); Sub-09's errors extended along the full
L--M opponent range with additional misidentification of magenta (c8).
The 2--12~nm range of L/M cone opsin polymorphisms in anomalous
trichromats~\cite{neitz2011, deeb2005} and the $\sim$8\% prevalence of
male CVD~\cite{bosten2019} locate these two cases within the
physiologically and epidemiologically relevant window.

Because only two symptomatic CVD cases were available, all inferences
below are framed as single-case analyses in the Crawford--Howell
tradition~\cite{crawford1998, schuett2023}; no group-level CVD claim is
made.

%-----------------------------------------------------------------------
\subsection{LORO preserved --- the filter precondition}
\label{sec:results:loro}

hV4 in CVD preserves categorical discrimination among the target hues ---
the minimum precondition for the corrective-filter programme. We
established this using leave-one-run-out (LORO) cross-validation of the
forward encoding model.

Classification accuracy in hV4 exceeded the 0.125 chance rate in every
individual (HC $= 0.68 \pm 0.11$, CVD $= 0.64 \pm 0.09$, mean $\pm$ SD
across 6 runs), and a between-group permutation test showed no HC--CVD
difference at any retinotopic ROI (hV4 $p = 0.668$; V1 $p = 0.542$;
V2 $p = 0.611$; 10{,}000 label permutations). The same pattern held under
SRM-aligned features. In line with previous work on CVD identification
performance at above-threshold contrast~\cite{bosten2019, boehm2014}, this
result confirms that the 8 scanner stimuli are individually distinguishable
at the cortical level --- the representational substrate on which a filter
could act is present. LORO therefore establishes the baseline applicability
of the corrective-filter programme in both CVD participants, and provides the
null against which the continuous-geometry phenotype
(\cref{sec:results:loco}) is contrasted.

%-----------------------------------------------------------------------
\subsection{LOCO vulnerability --- the filter target}
\label{sec:results:loco}

We next asked where the continuous hue geometry fails. Leave-one-colour-out
(LOCO) cross-validation trains the encoder on 7 of 8 hues and measures
prediction error for the held-out hue, yielding an 8-dimensional
vulnerability vector $\mathbf{v} \in \mathbb{R}^{8}$ per participant.

HC participants achieved above-chance hue interpolation in hV4
($\rho = 0.42 \pm 0.14$; group permutation $p = 0.044$), replicating
\citeA{brouwer2009}. Early visual cortex (V1, V2, V3) sustained reliable
discrimination (LORO) but failed continuous hue interpolation (LOCO not
significant): the hierarchical emergence of interpolation is confined to
hV4.

Both CVD participants fell near floor in hV4 (Sub-08 $\rho = 0.08$;
Sub-09 $\rho = 0.12$). HC--CVD effects were large across all retinotopic
ROIs (Hedges' $g = 1.61$, $1.85$, $1.34$ in V1, V2, hV4;
all $p_{\text{perm}} < 0.03$). The per-hue vulnerability pattern
differs between the two CVD participants in a way consistent with their
diagnostic axis: Sub-08's largest errors lie on c2/c3 (orange/yellow) and
c7 (purple); Sub-09's errors cluster on c5/c6 (cyan/blue) with magenta (c8)
as a further outlier (\cref{fig:main}b,c).
Relative to the HC reference range, the largest deviation for Sub-08 falls
at c7 (purple; Crawford--Howell $d = 2.40$, $p = 0.033$, one-tailed).
Both participants show pronounced deviation at c6 (blue):
Sub-08 $d = 2.13$, $p = 0.047$; Sub-09 $d = 2.15$, $p = 0.046$
(one-tailed single-case tests, uncorrected for 8 hues). Each
participant's vulnerability profile is an individually-structured target that
group-level classification metrics do not surface.

Behavioral concordance validates the LOCO phenotype as a functionally
interpretable metric. For the six hue-pair subsets probed in a same-day
just-noticeable-difference (JND) task, the hue pairs whose LOCO residuals
exceeded the individual HC reference range predicted above-JND confusion
with 100\% concordance. SRM $\Delta$RDM z-scores, measuring pairwise
representational geometry deviation, matched the same JND pairs at only
33\% concordance. The dissociation establishes that LOCO interpolation
capacity --- not pairwise geometric distance --- is the cortical quantity
that predicts behavioural colour discriminability.

\begin{figure*}[tb]
  \centering
  \includegraphics[width=\textwidth]{Figures/generated/fig1a_output.png}\\[0.6em]
  \includegraphics[width=0.85\textwidth]{Figures/generated/fig1_panels_bcd.pdf}
  \caption{\textbf{Subject-specific distortion fields recovered from
  hV4 hue representations.}
  \textbf{(a)} Pipeline. Procrustes-aligned fMRI amplitudes
  ($6 \times 8 \times V$) feed a forward encoder with $K{=}3$ half-wave
  rectified squared sinusoidal channels. LOCO cross-validation yields
  the vulnerability profile $\mathbf{v} \in \mathbb{R}^{8}$; this is
  fitted with the 2-component angular-dilation model ($\beta_s$, $\beta_c$),
  validated against the independent SRM $\Delta$RDM, and inverted to a
  per-hue stimulus-space correction filter.
  \textbf{(b)} Sub-08 (deutan): 2-component fit
  ($\beta_s = 38^\circ$, $\beta_c = -14^\circ$, $\rho = 0.881$,
  $p = 0.004$).
  \textbf{(c)} Sub-09 (protan): 2-component fit
  ($\beta_s = 6^\circ$, $\beta_c = -22^\circ$, $\rho = 0.690$,
  $p = 0.035$).
  Gray bars: observed LOCO Pearson $r$ per held-out hue
  (y-axis; $r = 0$ = chance); coloured line: model-predicted vulnerability
  under the 2-component model. $p$-values: $8! = 40{,}320$ exact
  permutations of colour labels on Spearman $\rho$.}
  \label{fig:main}
\end{figure*}

%-----------------------------------------------------------------------
\subsection{2-Component distortion model}
\label{sec:results:inference}

We fitted a two-component angular-dilation model in Stockman opponent-hue
space (Equation~\ref{eq:2comp}) to the hV4 LOCO vulnerability vector
$\mathbf{v}$: one component ($\beta_s$) weights distortion along the
S-cone opponent axis, and a second ($\beta_c$) along the CVD-subtype
confusion axis. Retinal-family comparators were evaluated in parallel
(Supplementary~S16 for full fits).

\paragraph{2-component model --- consistent parameterization structure across both subjects.}
The 2-component model reaches significance for both CVD participants with a
consistent model structure: Sub-08 ($\beta_s = 38^\circ$, $\beta_c =
-14^\circ$, $\rho = 0.881$, $p_{\text{perm}} = 0.004$) and Sub-09
($\beta_s = 6^\circ$, $\beta_c = -22^\circ$, $\rho = 0.690$,
$p_{\text{perm}} = 0.035$). The S-cone axis component ($\beta_s$) receives
non-zero weight in both cases, consistent with the behavioural evidence that
CVD perceptual errors concentrate near S-cone confusion
loci~\cite{emery2021}. The confusion-axis component ($\beta_c$) weights
the residual cortical structure not captured by the S-cone axis alone. The two components
together constitute a single parameterisation family applicable across
both subjects that admits an exact pre-image for all 8 hues at each
individual severity.

Retinal-family fits (Machado and R+C; Supplementary~S16) identify a
different best-fit model class per subject, preventing a common-axis
comparison of the two cases' distortion mechanisms --- an outcome that
motivates the independent selection criterion described below.

\begin{table}[tb]
  \caption{2-component distortion model fits at hV4 LOCO for both CVD
  participants. Grid search over $26 \times 51 = 1{,}326$ ($\beta_s$,
  $\beta_c$) combinations; significance by $8! = 40{,}320$ exact
  permutations of colour labels on Spearman $\rho$. Both participants
  reach significance ($p \leq 0.035$) with a consistent parameterisation
  structure across deutan and protan subtypes. Retinal-family
  comparators (Machado, R+C) in Supplementary~S16.}
  \label{tab:main}
  \centering
  \small
  \begin{tabular}{llccccc}
    \toprule
    Subject & Model & DOF &
    Param.\ 1 & Param.\ 2 &
    $\rho$ & $p_{\text{perm}}$ \\
    \midrule
    Sub-08 (deutan) & 2-component & 2 & $\beta_s = 38^\circ$ & $\beta_c = -14^\circ$ & 0.881 & \textbf{0.004} \\
    Sub-09 (protan) & 2-component & 2 & $\beta_s =  6^\circ$ & $\beta_c = -22^\circ$ & 0.690 & \textbf{0.035} \\
    \bottomrule
  \end{tabular}
\end{table}

\paragraph{Model selection: bijectivity as the deciding criterion.}
Despite comparable LOCO fit quality with retinal-family models
(Supplementary~S16), the 2-component model is selected as the primary
filter substrate on the criterion of pre-image bijectivity: for the
stimulus-space correction to be exact and unique for every displayed hue,
the forward distortion must be invertible. The 2-component angular dilation is a
bijective map on the opponent hue circle by construction, guaranteeing an
exact pre-image for all 8 hues at both severity levels. Machado
arc-compression at Sub-09's severity renders exact pre-images unavailable
for 4 of 8 hues (Supplementary~S16), violating injectivity.
The 2-component model avoids arc compression entirely: pre-image solutions
are exact for all 8 hues in both subjects, with numerical residuals
$< 0.001^\circ$.

%-----------------------------------------------------------------------
\subsection{SRM $\Delta$RDM --- independent geometric convergence}
\label{sec:results:srm}

Before proceeding to the corrective filter we asked whether the 2-component
distortion field also explains the pairwise representational geometry
independently of the LOCO criterion. This $\Delta$RDM validation is
statistically independent of the LOCO fitting procedure.

Voxel-space crossnobis distances among the 8 hues were SRM-projected into a
shared low-dimensional space trained on HC participants only (7~HC;
$K = \{4, 4, 3, 3\}$ for V1, V2, V3, hV4, respectively). For each CVD
participant we computed
$\Delta\text{RDM} = \text{RDM}_{\text{CVD}} - \overline{\text{RDM}}_{\text{HC,LOO}}$
and tested whether the 2-component distortion parameters, when
forward-simulated, predict the observed $\Delta$RDM
pattern~\cite{schuett2023, chen2015, bannert2025}.

For Sub-09, the 2-component model reaches significance under both the
functional LOCO criterion and the independent geometric $\Delta$RDM
criterion ($p = 0.007$; 40{,}320 cosine-similarity permutations on
crossnobis distances), constituting stronger case-level evidence than
either criterion alone. An independent cross-model convergence check
using R+C parameters also reaches $p = 0.026$ for Sub-09 (Supplementary~S16).
For Sub-08, neither criterion reaches significance in $\Delta$RDM at any
ROI ($p > 0.05$). The Sub-08 case therefore rests on LOCO significance
and pre-image bijectivity alone.

$\Delta$RDM is reported as an \emph{auxiliary convergence check}, not as a
primary fitting criterion. SRM-alignment partially absorbs cone-shift signal
into the shared basis in these data. Retinal-family optimisation against
$\Delta$RDM fails for both CVD cases (Supplementary~S16), confirming that
LOCO is the more informative phenotype and $\Delta$RDM is best used as a
cross-check on the LOCO-derived parameters.

%-----------------------------------------------------------------------
\subsection{Pre-image inversion --- individual correction filters}
\label{sec:results:filter}

Given the validated 2-component distortion field, the per-hue correction
filter is the stimulus-space pre-image of the distortion: for each displayed
hue $\theta_k$, the corrected stimulus $\tilde{\theta}_k$ satisfies
$T(\tilde{\theta}_k) = \theta_k$, where $T$ is the forward angular dilation
(\cref{sec:methods:filter}). Pre-images were computed numerically by
dense-grid initialisation followed by Brent's method refinement
(\cref{sec:methods:filter}).

\paragraph{Bijectivity and correction vectors.}
The 2-component filter recovers all 8 hues exactly for Sub-08 (mean
$|\delta| = 46.3^\circ$, max $= 104.2^\circ$, residual $< 0.001^\circ$)
and all 8 for Sub-09 (mean $|\delta| = 20.1^\circ$, max $= 48.1^\circ$,
residual $< 0.001^\circ$). By contrast, Machado arc-compression at
Sub-09's severity yields exact pre-images for only 4 of 8 hues
(Supplementary~S16). The collapsed hues
fall within Sub-09's LOCO vulnerability cluster (\cref{fig:main}c),
confirming that arc-compression eliminates the correction exactly where
Sub-09 requires it most.

\paragraph{Detection agrees; correction diverges.}
All models identify the same two participants as requiring correction.
The 2-component and retinal-family correction vectors for Sub-08,
however, diverge markedly at comparable LOCO fit (Supplementary~S16).
Detection converges; correction directions diverge. This divergence
yields the study's primary behavioural prediction. A 2AFC task comparing
2-component and retinal-family filter conditions will identify whether
cortical or retinal distortion is the operative corrector.

\paragraph{Individual filter divergence.}
The correction vectors for Sub-08 and Sub-09 are individually structured
and divergent: the two subjects' filter profiles have cosine similarity
$< 0$ across hue pairs, confirming that a single population-average
correction applied to both would fail at least one individual. The
idiosyncratic structure of each filter is a direct consequence of the
individually-fitted ($\beta_s$, $\beta_c$) parameters and is not captured
by any retinal group-average model.

\begin{figure*}[tb]
  \centering
  \includegraphics[width=0.95\textwidth]{Figures/generated/fig2_output.png}
  \caption{\textbf{2-component pre-image inversion yields exact corrections
  at both severity levels.}
  \textbf{(a)} Sub-08 (deutan, 2-component: $\beta_s = 38^\circ$,
  $\beta_c = -14^\circ$): 8/8 exact pre-images (gray arrows;
  $< 0.001^\circ$ residual; mean $|\delta| = 46.3^\circ$).
  \textbf{(b)} Sub-09 (protan, 2-component: $\beta_s = 6^\circ$,
  $\beta_c = -22^\circ$): 8/8 exact pre-images (mean
  $|\delta| = 20.1^\circ$). Dashed outlines: hues where Machado
  arc-compression ($\Delta\lambda = 13.5$~nm) collapses three distinct hues
  to a single percept, rendering exact Machado pre-images unavailable for 4
  of 8 hues. The 2-component filter avoids arc compression entirely.}
  \label{fig:feasibility}
\end{figure*}

\paragraph{Case-level specificity.}
Case-level specificity is supported by three converging grounds,
independent of the elevated false-positive rate under grid-search
permutation testing alone. Applied to the 7 HC participants as CVD
candidates, the pipeline reaches nominal $p < 0.05$ in 15 of 21
model--subject tests (\cref{app:hc_specificity}), driven by
rank-based-loss flexibility over 8 data points and regression to the
mean on per-HC baseline $\rho$. This rate confirms that permutation
significance alone does not distinguish CVD from HC. The three grounds
are: (i) for
Sub-09, the LOCO criterion and the independent $\Delta$RDM criterion both
reach significance under the 2-component parameterisation, and no HC
achieves cross-criterion significance in both; (ii) the Sub-09 retinal-family
fit magnitude falls within the physiologically interpretable range for
moderate--severe protanomaly (Supplementary~S16;~\citeNP{neitz2011}),
whereas HC-significant fits distribute across the grid with no physiological
anchor; (iii) the 2-component S-cone axis structure is independently
motivated by behavioural evidence that CVD errors concentrate near S-cone
confusion loci~\cite{emery2021}, providing a physiological grounding for the
model form independent of the fit quality. Together these three anchors
provide case-level specificity. An explicit behavioural test of the
pre-image filters~(\cref{sec:results:behavioralfilter}) is the preregistered
falsifier for the cortical-level account.

% --- Behavioural validation subsection (Phase 3 pending) ---
% \subsection{Behavioural validation of the pre-image filters}
% \label{sec:results:behavioralfilter}
% Placeholder — JND experiment comparing uncorrected display, R+C filter,
% and 2-component filter for Sub-08 and Sub-09. Preregistered prediction:
% along confusion-axis hue pairs from each subject's LOCO profile, only the
% cortical 2-component filter will significantly reduce JND relative to
% uncorrected; the Machado filter will fail to improve or worsen performance
% for Sub-09 (arc compression) and show opposing correction directions for
% Sub-08 (cosine similarity −0.18).

%================================================================
% Cross-reference targets used:
%   fig:main          — Figure 1 (pipeline + LOCO profiles)
%   fig:feasibility   — Figure 2 (pre-image / 2-comp vs Machado)
%   tab:main          — Table 1 (model fits per CVD participant)
%   tab:full          — Supplementary (cross-ROI × cross-criterion)
%   app:hc_specificity — Supplementary (HC FPR calibration)
%   sec:methods:filter — Methods §Filter design
%================================================================
